\section{REVIEW OF EXISTING ALTERNATIVES}
Staff from the International School Manila and outside reviewers currently depend on centralized attestation, registrar portals, and provider immutability for skew authenticity.The current International School Manila's staff and outside evaluators rely on centralized attestation, registrar portals, and provider immutability as proof of skew authenticity. Access logs and versioning are implemented for files.In certain higher education institutions (HEIs) in the Philippines, Formal checks are done via CHED eCAV that standardizes the requests for the checks in a platform-bound and location-addressed manner. Verifiers should only believe the metadata from the portal and should not try to validate a content bound proof from the bytes of the file \citep{CommissionOnHigherEducationND}. Pilots, such as DPWH Integrity Chain and GovChain, show national enthusiasm for tamper evident public records but only address public sector records and cannot oversee the range of academic reports, each with its own separate and independent third party authentication of the notoriously fraught report-by-byte level. \citep{DPWH2025,BetterGovernmentPhilippinesND}.

Education is based on verifiable credentials as the strongest resources. EduCredPH forwarded a concept of a permissioned blockchain network for educational credential verification in the Philippine context where HED and third party verifiers can issue credentials and verified on a common Hyperledger Fabric network, all while learners control and consult their credentials \citep{EduCredPH2024}.On OpenCerts, this is extended across institutions by each institution minting a signed payload which can be signed by any verifier without communication with the issuer. The technical annex of the SkillsFuture Transformation provides an understanding of the data model and the verification flow in easy-to-understand operational steps and the work with the ones NTU has demonstrated during the institution's adoption of the SkillsFuture Transformation is illustrated in a daily fashion \citep{GovTech2019,SkillsFuture2018,NTUND}. End to end issuance, revocation, and verifier experience with research and university application demonstrates feasibility in ASEAN and even worldwide areas. All such efforts continue to be credential-centric and fail to offer content addressed permanence for unstructured academic reports on scale.\citep{Naji2023,MohebiAghdam2024}.

VERACISM is unique in taking the credential verification ideology to each completed report. When ingesting a file is scanned and validated, and then hashed to retrieve its content identifier (SHA256) which differs if any of the file's bytes are altered. Files are stored on decentralized infrastructure for which deal evidence can be seen using public explorers for durability, and for public corroboration. Auditors, educators and students can rehash a copied file and compare the rehashed identifier to the recorded identifier, without any ISM credentials, and confirm storage evidence. This design is putting cryptographic proofs under the current portals. eCAV can continue to be a national front door for policy and records processing, while VERACISM provides content addressed evidence and independent verification for any report artifact. Local solutions like GovChain and Integrity Chain prove that this is feasible on public records that need to be tamper evident. VERACISM extends this concept to the academic setting by supporting role based access, tenant aware metadata, and immutable audit events. In the end, VERACISM extends beyond the concept of focusing on credentials or centralizing trust; it's the tangible and independently verifiable evidence offered in every report.\citep{CommissionOnHigherEducationND,BetterGovernmentPhilippinesND,DPWH2025,GovTech2019,SkillsFuture2018,NTUND,EduCredPH2024,Naji2023,MohebiAghdam2024}.

\begin{center}
\captionsetup{type=table}

\smalltable
\scriptsize
\begin{tabular}{|>{\raggedright\arraybackslash}p{2.6cm}|
                >{\raggedright\arraybackslash}p{1.8cm}|
                >{\raggedright\arraybackslash}p{2.6cm}|
                >{\raggedright\arraybackslash}p{2.6cm}|
                >{\raggedright\arraybackslash}p{2.6cm}|}
\hline
\textbf{System (License/Type)} &
\textbf{Platform} &
\textbf{Brief Description} &
\textbf{Basic Functionalities} &
\textbf{Relation to the Proposed System} \\
\hline\hline

\textbf{VERACISM} &
.NET 8 API, MSSQL, IPFS, Filecoin &
Content addressed storage and verification for ISM reports. &
Scan and validate files, hash with SHA 256, store on decentralized storage, log actions, enforce role based access. &
-- \\
\hline

CHED eCAV (Government portal) &
Web portal &
National portal for electronic verification of academic records. &
Receive verification requests, check records, issue electronic certifications. &
External authority; VERACISM can add file level proof for ISM records used in eCAV. \\
\hline

DPWH Integrity Chain and GovChain (Government pilots) &
Public blockchain platforms &
Pilots for transparent and tamper evident government project data. &
Record project data on a blockchain and expose entries for public viewing and checking. &
Demonstrate feasibility of tamper evident records but focus on government datasets, not ISM reports. \\
\hline

EduCredPH (Philippine credential network) &
Permissioned blockchain (Hyperledger Fabric) &
Credential network for Philippine higher education institutions. &
Issue, store, and verify academic credentials across participating institutions. &
Addresses credential fraud and multi institution verification, but not unstructured ISM reports or malware scanning. \\
\hline

OpenCerts (Singapore) (National system) &
Public verification system &
Singapore platform for digital education and training certificates. &
Issue structured certificate files and support public verification without contacting issuers. &
Mature credential model but tied to its ecosystem and not designed as general ISM report storage. \\
\hline

Research systems: UTM BADVES and DIAR (Research) &
University prototypes &
Blockchain based academic credential issuance and checking. &
Model issuing, revoking, and verifying awards using distributed ledger technology. &
Provide design patterns and feasibility evidence but are prototypes, not ready for ISM production with diverse reports. \\
\hline

\end{tabular}
\caption{Summary of Existing Alternatives}
\label{table:existing-alternatives-summary}
\end{center}


\begin{center}
	\captionsetup{type=table}
	
    \smalltable \scriptsize
	\begin{tabular}{|>{\raggedright\arraybackslash}p{3.0cm}| >{\centering\arraybackslash}p{1.7cm}|
	>{\centering\arraybackslash}p{1.7cm}| >{\centering\arraybackslash}p{1.7cm}| >{\centering\arraybackslash}p{1.7cm}|
	>{\centering\arraybackslash}p{1.7cm}| >{\centering\arraybackslash}p{1.9cm}|}
		\hline
		\textbf{System}                 & \textbf{Content-\ addressed storage for all reports} & \textbf{Decentralized storage for reports} & \textbf{Malware scanning at ingest} & \textbf{Credential-\ focused model} & \textbf{Government / public - \ records focus} & \textbf{Internal ISM storage and checking layer} \\
		\hline
		\hline
		\textbf{VERACISM}               & Yes                                                 & Yes                                        & Yes                                 & No                                 & No                                          & Yes                                              \\
		\hline
		CHED eCAV                       & No                                                  & No                                         & No                                  & Yes                                & No                                          & No                                               \\
		\hline
		DPWH Integrity Chain / GovChain & No                                                  & Yes                                        & No                                  & No                                 & Yes                                         & No                                               \\
		\hline
		EduCredPH                       & No                                                  & No                                         & No                                  & Yes                                & No                                          & No                                               \\
		\hline
		OpenCerts (Singapore)           & No                                                  & No                                         & No                                  & Yes                                & No                                          & No                                               \\
		\hline
		UTM BADVES / DIAR               & No                                                  & No                                         & No                                  & Yes                                & No                                          & No                                               \\
		\hline
	\end{tabular}
    \caption{Feature Mapping of Existing Alternatives}
	\label{table:existing-alternatives-feature-mapping} 
\end{center}



\subsection{Summary of Existing Alternatives}

The systems in Table~\ref{table:existing-alternatives-summary} each solve a part of the trust and verification problem, but they are designed for their own scope. Together they strengthen credentials and public sector records, yet they still leave ISM without a single, content based pipeline for all finalized reports under local control.

CHED eCAV is a national gateway to the verification of academic records, however, it only receives requests and provides responses and does not access the files and does not encroach on the protection provided by ISM on how it stores them. DPWH Integrity Chain and GovChain demonstrate that the government project-related information can be placed on a blockchain, but it is configured to student reports. EduCredPH provides higher education schools with a common means of issuing and verifying credentials, whereas OpenCerts provides Singapore with the same but nationwide, both remain limited to formal certificates. UTM BADVES and DIAR university prototypes investigate the way a university can produce and revoke credentials using distributed ledger technology, but these are not yet products, only research projects. These systems combined aid in external verification, credential fraud, but each of them cannot serve as the internal storage and checking layer of ISM to all report files, and thus cannot substitute VERACISM.

The biggest gap in the existing systems is that they look either at credentials alone or at government project data, so none of them gives ISM lasting protection for every report together with malware checks and audit trails that follow ISM roles and students. VERACISM is worth building because it pulls those missing pieces into one flow that staff already use. When a report is uploaded, the system first runs antivirus and YARA scans, then checks the file type and size, computes a hash, and saves the content on decentralized storage while writing every step into an audit log that cannot be changed. It treats each finalized report much like a diploma, giving it a content identifier that is tied to ISM metadata and access rules. Afterwards any authorized person can retrieve the file, recalculate the hash, and compare the result with the one stored to assure that the report hasn't been changed and is still available in storage. This is a change to the habit of trust around ISM records and makes it visible and verifiable.
